% ============================================================================
%  DATOS DEL TFL — completar esto y no hace falta tocar nada mas.
%  Estos datos se usan en la portada, en los metadatos del PDF y en el
%  encabezado de las paginas.
% ============================================================================

% --- Titulo y autor/a ---------------------------------------------------
\newcommand{\TFLTitulo}{Geología de las nacientes del río Malargüe, provincia de Mendoza}
\newcommand{\TFLSubtitulo}{}                    % opcional, se puede dejar vacio
\newcommand{\TFLAutor}{Ricardo Tapia}
\renewcommand{\TFLApellido}{Tapia}              % aparece en el encabezado de cada pagina

% --- Direccion -----------------------------------------------------------
\newcommand{\TFLDirectores}{Dr. Bruno Díaz}
\newcommand{\TFLCodirectores}{}                 % opcional, se puede dejar vacio

% --- Carrera e institucion ------------------------------------------------
\newcommand{\TFLCarrera}{Licenciatura en Ciencias Geológicas}
\newcommand{\TFLInstitucion}{Universidad de Buenos Aires}
\newcommand{\TFLFacultad}{Facultad de Ciencias Exactas y Naturales}
\newcommand{\TFLDepartamento}{Departamento de Ciencias Geológicas}

% --- Lugar y fecha de entrega ---------------------------------------------
\newcommand{\TFLLugar}{Ciudad Autónoma de Buenos Aires}
% \today se actualiza automaticamente cada vez que se compila el documento.
% Para fijar una fecha manual, reemplazar \today por el texto deseado, por
% ejemplo: \newcommand{\TFLFecha}{Abril de 2027}
\newcommand{\TFLFecha}{\today}

% --- Metadatos del PDF (no se ven en la portada, van "adentro" del archivo) --
\newcommand{\TFLPalabrasClave}{geología, río Malargüe, Mendoza, TFL, FCEN-UBA}
\renewcommand{\TFLIdioma}{es-AR}
\newcommand{\TFLResumenCorto}{Trabajo Final de Licenciatura sobre la geología de las nacientes del río Malargüe, provincia de Mendoza.}

% --- Logos e imagenes de portada ------------------------------------------
% Deben ser archivos de imagen (png/jpg) dentro de logos/ o figuras/.
% TFLLogoDerecho va arriba a la izquierda, dentro de la banda de color;
% TFLLogoIzquierdo va como sello abajo a la derecha (los nombres son
% historicos y no describen la posicion final en la portada).
\newcommand{\TFLLogoIzquierdo}{logos/Logo-fcenuba.png}
\newcommand{\TFLLogoDerecho}{logos/Logo_Geologia_oficial.png}

% Imagen opcional para reemplazar la banda de color de la portada por una
% foto. No se usa por defecto (ver \TFLPortada en main.tex).
\newcommand{\TFLImagenPortada}{figuras/IMG_20160223_152527650.jpg}
